\chapter{Conclusioni e sviluppi futuri}
\label{chap:conclusioni}

La tesi ha confrontato il sensore PIR con due moduli radar mmWave a 24~GHz per il rilevamento della presenza nei nodi del progetto DIPME. Lo studio dei sensori e della loro documentazione ha chiarito il principio di funzionamento e i dati che ciascuno fornisce, dal singolo bit del PIR alla distanza e all'energia per gate del radar. Per il confronto è stata realizzata una piattaforma di acquisizione su ESP32, con un'interfaccia web servita dal nodo stesso che mostra i dati in tempo reale e li esporta in CSV.

La campagna sperimentale ha confermato che il PIR, per principio, non vede una persona immobile, nemmeno sotto il banco. A decidere se il PIR vede è il movimento, non la distanza. L'\ldb\ invece rileva la persona immobile in tutti gli scenari, attraversa legno, vetro, cartongesso e plastica, che bloccano il PIR, e non ha prodotto falsi positivi a stanza vuota. L'\ldc\ funziona nello stesso scenario, ma non localizza la persona ferma e richiede una taratura per installazione, ed è quindi meno adatto. Sulle grandezze del radar è stato poi costruito un indice di vitalità, calcolato a bordo del nodo, che distingue la persona immobile da quella in movimento, con soglie che dipendono dalla geometria in cui sono tarate.

Il radar consuma però circa mille volte più del PIR e, sempre acceso, esaurisce la batteria in meno di un giorno. La risposta alla domanda di partenza è quindi che il radar rileva la persona con maggiore affidabilità, ma deve affiancare il PIR, non sostituirlo. In tempo di pace il nodo resta in deep-sleep e il PIR lo sveglia quando rileva una presenza, come già avviene nel progetto~\cite{callisto2026dipme}. In emergenza si accende il radar, che rileva anche la persona ferma, e alternandolo con un ciclo di lavoro l'autonomia supera i tre giorni. La raccomandazione per il progetto è di aggiungere al nodo un modulo della classe dell'\ldb\ secondo questa architettura, con una taratura all'installazione sotto l'arredo.

\section{Sviluppi futuri}
\label{sec:conclusioni-sviluppi}

Restano tuttavia diverse aree da esplorare. La prima è l'integrazione con la rete LoRa. Nel regime di emergenza il bit di presenza decide la cadenza dei messaggi del nodo~\cite{callisto2026dipme}, e sostituirlo con la presenza del radar è l'innesto più immediato, mentre l'indice di vitalità, che occupa un byte, renderebbe i messaggi più informativi.

Un'altra riguarda la prova nel dimostratore reale, dove restano da verificare i falsi positivi con arredi metallici, il comportamento con più persone, il montaggio rispetto alla lamiera forata del banco, la separazione fra banchi adiacenti e il rilevamento immediato della persona quando il radar viene riacceso.

Infine, una procedura di auto-taratura all'installazione adatterebbe le soglie dell'indice a ogni banco, e un confronto a tre con il sensore UWB previsto nelle slide del progetto~\cite{slideDipme} porterebbe a una misura diretta il confronto oggi basato sulla letteratura.

In sintesi, il radar mmWave colma il limite del PIR nello scenario per cui il progetto DIPME è nato, e la piattaforma realizzata è pronta a portare il confronto nelle aule reali.
